\begin{abstract}
许多张量算子将一个归约的最终结果用作另一个归约的上下文。
这种依赖阻碍了常规循环融合，使单遍算法仍需专家推导和实现。
融合需要保留足够的信息，以便在上下文变化时修正部分结果，并通过合并操作（join）完成这些修正，而不重新读取输入。
本文提出 \sys，借鉴自动化分治并行的状态合成与合并构造方法，实现依赖归约的循环融合。
代数分析得到能够独立合并、并确定所需结果的参考状态。
直接执行这一表示可能展开源程序的中心化和缩放，从而引入相消或溢出。
因此，\sys 以源归约结果为基础，结合分析上下文修正时发现的补充量，构造维护状态。
语法引导合成组合这些量，使维护状态能够恢复参考状态。
两种表示之间经验证的转换导出在实数语义下正确的合并操作。
随后将合并重写为直接修正，尽可能保留源程序的中心化和缩放，以尽量保持其数值精度。
当每个分区只有一个元素时，这些状态可以在一次融合遍历中计算并合并；同一合并操作也支持并行归约树。
TVM 实现在选取的 \nShowcase 个算子中，为 \nShowFused 个推导出单遍归约器；实验评估覆盖率、发现的状态、推导开销和数值行为。
\end{abstract}
